% 以 draft/D题/main.tex 为准；这里只抽出与论文内容无关的导言设置。
% 从 draft/D题 提取的统一表格样式；供所有三线表复用。
\newcolumntype{Y}[1]{>{\centering\arraybackslash}m{#1}}
\newcommand{\papertablestyle}{%
  \normalfont\songti\zihao{-4}%
  \linespread{1.15}\selectfont%
  \renewcommand{\arraystretch}{1}%
  \setlength{\tabcolsep}{5pt}%
  \setlength{\heavyrulewidth}{1pt}%
  \setlength{\lightrulewidth}{0.5pt}%
  \setlength{\aboverulesep}{3pt}%
  \setlength{\belowrulesep}{3pt}%
  \captionsetup{font={normalfont,song,minusfour},position=top,skip=4pt}%
}

\usepackage{fvextra,tcolorbox,pdfpages,eso-pic}
\tcbuselibrary{breakable}
\clubpenalty=10000
\widowpenalty=10000
\displaywidowpenalty=10000
\setCJKfamilyfont{abstractlisu}[Path=fonts/]{LiSu.ttf}
\geometry{footskip=13.82777pt}
\setcounter{tocdepth}{3}

\makeatletter
\def\@oddfoot{\hfill{\normalfont\fontsize{9bp}{11bp}\selectfont\thepage}\hfill}
\let\@evenfoot\@oddfoot
\newcommand\papertocline[4]{%
  \par\begingroup\normalfont\songti\fontsize{12bp}{18.6bp}\selectfont
  \leftskip=#1\relax\rightskip=12bp\parfillskip=-\rightskip
  \parindent=0pt\interlinepenalty=10000
  \@tempdima=#2\relax\advance\leftskip\@tempdima
  \leavevmode\hskip-\@tempdima{#3}\nobreak\hskip3bp
  \leaders\hbox to3bp{\hss.\hss}\hfill\nobreak
  \hb@xt@12bp{\hfil\normalfont\fontsize{12bp}{18.6bp}\selectfont#4}\par
  \endgroup}
\renewcommand\l@section[2]{\papertocline{0bp}{30bp}{#1}{#2}}
\renewcommand\l@subsection[2]{\papertocline{12bp}{21bp}{#1}{#2}}
\renewcommand\l@subsubsection[2]{\papertocline{24bp}{30bp}{#1}{#2}}
\renewcommand\tableofcontents{%
  \begingroup\normalfont\songti\fontsize{12bp}{18.6bp}\selectfont
  \setlength{\parskip}{0pt}\setlength{\parindent}{0pt}
  \vspace*{0bp}
  {\centering\fontsize{16.2bp}{20bp}\selectfont 目录\par}
  \vspace{2bp}\@starttoc{toc}\endgroup}
\makeatother
